\section{Introduction}
Cancer remains one of the leading causes of death worldwide and therefore
represents a major challenge for modern medicine. A crucial factor for
successful treatment is the early and reliable detection of cancerous tissue.
In many diagnostic workflows, tissue samples obtained through biopsies or
surgical procedures are examined under a microscope by trained pathologists.
This process, known as histopathological analysis, allows medical experts to
identify structural abnormalities and cellular patterns that are associated with
cancer. However, the increasing amount of medical imaging data makes manual
analysis both time-consuming and expensive. Furthermore, the diagnostic process
heavily depends on the experience and concentration of human experts, which can
lead to inter-observer variability and diagnostic inconsistencies.

To address these challenges, machine learning and deep learning approaches have
become increasingly important in medical image analysis. In particular,
convolutional neural networks (CNNs) have shown strong performance in image
classification tasks due to their ability to automatically learn hierarchical
visual features directly from image data. Architectures such as ResNet-50 are
therefore frequently used for histopathological image classification. Instead of
relying entirely on handcrafted image features, CNN-based systems can learn
discriminative tissue characteristics directly from labeled training data. This
enables the development of computer-aided diagnostic systems that can support
pathologists by accelerating the analysis process and highlighting suspicious
tissue regions.

Despite the promising performance of deep learning approaches, several
challenges remain. Histopathological images are highly complex and often exhibit
substantial visual variability caused by staining differences, tissue
deformations, and imaging artifacts. In addition, obtaining accurately labeled
medical data is difficult because annotations require expert knowledge from
trained pathologists. Another important challenge is the high cost of
classification errors. In particular, false negatives, where cancerous tissue is
classified as non-cancerous, can delay diagnosis and treatment and therefore
have severe medical consequences. Consequently, optimizing metrics related to
cancer detection performance, such as recall for the cancer class, is often more
important than maximizing overall accuracy alone.

\subsection{Dataset}
To investigate these challenges, this work focuses on the classification of
histopathological image patches using the
PathMNIST\footnote{https://medmnist.com/} dataset. The dataset contains
microscopic image patches extracted from colorectal cancer histology slides and
categorized into multiple tissue classes, including tumor tissue and several
non-cancerous tissue types.

The complete dataset consists of 107,180 samples. These samples are divided into
89,996 training samples, 10,004 validation samples, and 7,180 test samples. The
training set is used to fit the model parameters, the validation set is used
during model selection and hyperparameter tuning, and the test set is reserved
for the final evaluation. Each sample is available in two resolutions: a compact
$28 \times 28$ pixel version and a higher-resolution $224 \times 224$ pixel
version. The $28 \times 28$ version allows efficient experimentation with
lightweight models, while the $224 \times 224$ version preserves more spatial
detail and is therefore more suitable for deeper convolutional neural networks
such as ResNet-based architectures.

PathMNIST contains nine tissue classes: adipose, background, debris,
lymphocytes, mucus, smooth muscle, normal colon mucosa, cancer-associated
stroma, and colorectal adenocarcinoma epithelium.
Figure~\ref{fig:pathmnist-class-examples} shows one representative example image
for each class. These classes cover both normal and pathological tissue
structures. Adipose corresponds to fat tissue, while background denotes image
regions without relevant tissue content. Debris describes unstructured tissue
fragments or cellular remains. Lymphocytes are immune cells that may occur in
the tumor microenvironment or surrounding tissue. Mucus refers to mucinous
material, which can be present in colorectal tissue and certain tumor regions.
Smooth muscle represents muscular tissue from the intestinal wall. Normal colon
mucosa corresponds to healthy colon epithelial tissue with regular glandular
structures. Cancer-associated stroma describes connective tissue associated with
tumor growth and the tumor microenvironment. Finally, colorectal adenocarcinoma
epithelium represents malignant epithelial tumor tissue and is therefore the
main class of interest in this work.

\begin{figure}[ht]
    \centering
    \includesvg[inkscape=false,inkscapelatex=false,width=0.80\textwidth]{figures/pathmnist_class_examples}
    \caption{Representative image examples for the nine PathMNIST classes.}
    \label{fig:pathmnist-class-examples}
\end{figure}


The class distribution across the training, validation, and test splits is shown
in Figure~\ref{fig:pathmnist-class-distribution}. The distribution is not
perfectly uniform across all classes. In the training and validation sets, the
largest class fractions are observed for smooth muscle and colorectal
adenocarcinoma epithelium, while classes such as normal colon mucosa, mucus, and
lymphocytes occur with slightly smaller but still substantial proportions. In
contrast, the test split shows stronger deviations for some classes: adipose,
mucus, and colorectal adenocarcinoma epithelium have comparatively higher
fractions, whereas debris and cancer-associated stroma occur less frequently.
This is relevant for model evaluation because differences between the split
distributions can influence the measured performance on the test set.

\begin{figure}[ht]
    \centering
    \includesvg[inkscape=false,inkscapelatex=false,width=\textwidth]{figures/pathmnist_class_distribution}
    \caption{Class distribution of the PathMNIST dataset across the training, validation, and test splits.}
    \label{fig:pathmnist-class-distribution}
\end{figure}


The distribution shown in Figure~\ref{fig:pathmnist-class-distribution} also
motivates the choice of evaluation metrics. Since the aim of this work is to
reliably identify cancer-related tissue, overall accuracy alone is not
sufficient. A model may achieve high accuracy while still missing a relevant
number of cancerous samples, especially if errors are unevenly distributed
across classes. In this work, the primary cancer class is
\textit{colorectal adenocarcinoma epithelium}, which represents malignant tumor
tissue. In addition, \textit{cancer-associated stroma} is considered a relevant
secondary class because it describes tissue associated with the tumor
microenvironment. For this reason, special attention is given to the
classification performance for these cancer-relevant classes.

Instead of focusing primarily on accuracy, this work emphasizes the $F_2$ score
for the relevant classes. The $F_2$ score belongs to the $F_\beta$ family of
metrics, which combine precision and recall into a single evaluation measure and
is defined as follows:
$$ F_\beta = (1 + \beta^2) \cdot \frac{\text{precision} \cdot \text{recall}}{(\beta^2 \cdot \text{precision}) + \text{recall}} $$
Therefore, $F_2$ weights recall more strongly than precision, which is
particularly important in a medical context because false negatives, where
cancerous tissue is classified as non-cancerous, can delay diagnosis and
treatment. At the same time, the metric still incorporates precision, preventing
the classifier from trivially assigning all samples to the cancer class.
Consequently, the dataset is suitable for evaluating whether a classifier can
not only perform well globally, but also reliably detect cancer-related tissue
under a recall-oriented evaluation setting.

% TODO: Which classes should be important to detect right? I mean there are
% multiple classes related to cancer.
